\chapter*{引言}
\addcontentsline{toc}{chapter}{引言}
公众对艺人的认识往往来自多种媒介。同一人物在杂志中以视觉造型出现，在影视中通过角色进入故事，在社交平台中又以较短的图文或视频参与交流。将这些内容放在一起观察，可以讨论人物形象如何随媒介语境改变。

本文选择佐佐木希作为示例对象，围绕“不同媒介突出哪些人物特征”展开简要分析。自我呈现理论提供了理解表达场景与形象关系的一个视角，但不能据此推断人物的真实动机。\footnote{参见 Goffman E. The Presentation of Self in Everyday Life[M]. New York: Anchor Books, 1959.}
正文只建立比较框架，不涉及私人生活，也不把模拟材料写成实测结论。
